\section{Experiment A: Static Rostering}
\label{sec:exp-static}

% 证据：实验 5、实验 5 续；报告 实验5_A/B/C（+ T10 一句）
% 表 3：规则来源主表；表 4：最终方法三种设置；图 5：成本—代理分曲线；图 6：分组柱图
% 守住：
%  - 不写"LLM 规则优于 GP / 神经网络"
%  - 不写"在所有实例上优于 dec5"（E 组只是无显著差别）
%  - 不写"K0.5 是最优密度"（除非 T10 给出中间最优点）
%  - 方法成绩用各设计运行的平均，不用最好一次；两两比较报 Holm 校正后的 p
%  - 每个表注明参与归一化的臂集合

\subsection{Setup}
\todo{Train/test split (15/47), MOEA, budgets, statistics (Holm-corrected pairwise tests).}

\subsection{Rule Sources}
\todo{Three tiers; Table 3.}

\subsection{Final Method}
\todo{K0.5 + ES16/GP; Table 4; E/H/M groups (Fig.~6).}

\subsection{Stability, Interpretability and Design Cost}
\todo{Stability and H13 robustness diagnosis; GP converges to staggered maintenance; cost--proxy curve (Fig.~5).}
